
\section{Conclusions}
\label{sec:conclusions}

Moving from one patch to a network did not, by itself, change how augmentative releases of \textit{T.~radiata} should be judged. When dispersal was free, the number of parasitoids released set the pest load, uniform and targeted releases were comparable per parasitoid, and the graph was a minor modifier whose effect ran mainly through the degree of the outbreak patch. The linear stability of the synchronous state ordered the networks by their spatial coherence, but the level of coherence reached over an outbreak was set mainly by transient dynamics.

Transit mortality is what made network structure matter. It produced an optimal parasitoid dispersal rate for targeted releases, it broke the equivalence between the star and the complete graph, and it reversed the benefit of a well-connected outbreak patch. The optimum follows a persistence threshold for the resident parasitoid that depends on the mean degree and the mean edge mortality rather than on the outbreak patch alone. Because it marks the point where natural control fails, a high efficiency per released parasitoid can coexist with a high pest load, and programmes should report both.

For \textit{D.~citri}, these results favour area-wide releases when parasitoid mobility is low, in absolute terms or relative to a psyllid introduced at one point, and give no support to targeting the best-connected grove as a general rule. Two measurements would decide which regime applies to citrus landscapes, the relative dispersal rates of psyllid and parasitoid between groves and the mortality of \textit{T.~radiata} while moving between them. A calibrated, stochastic version of the network model would be needed before these conclusions are translated into release schedules.
